% Insert after the preserved as-is observation tail. Keep its existing section.
\hyphenrules{english}
\subsection{Mission requirements and effectiveness}
\label{subsection:framework-mission-requirements}
\label{subsection:framework-mission-needs}
\label{subsection:framework-moes}
\label{subsection:eps-reference-lifecycle}

These proposed enterprise requirements state EPS outcomes, while the MOEs express stakeholder value, following OOSEM \citep[Sec.~17.3.2]{friedenthal2014practical}. The scope covers NASA Pre-Phase A, A, B and adapted C for modeled EPS hardware and host-executed control/interface software. It ends at model/software integration readiness before Phase D. Physical development, procurement, hardware operation and testing, PIL/HIL, target qualification and flight readiness are excluded. Hardware maturity relies on applicable existing evidence, with gaps recorded. NASA supplies reference phases and review purposes without a certification claim \citep[Secs.~2.2--2.4 and~3.3--3.6]{nasa2016systems}. A prepared review package is an entrance product. Acceptance requires the designated authority's decision against agreed success criteria \citep[Sec.~6.7 and Table~6.7-1]{nasa2016systems}.

% TODO: Agree measurement targets, acceptable errors and review thresholds with stakeholders.
% TODO: Review downstream scenario and requirement traces for the condensed MR statements.
% TODO: Refresh source locators, parser snapshots and publication checks after compression.
% Preserve historical review evidence and approved/protected source hashes.

{\small
\setlength{\tabcolsep}{4pt}
\renewcommand{\arraystretch}{1.08}
\begin{longtable}{@{}>{\raggedright\arraybackslash}p{0.09\linewidth}>{\raggedright\arraybackslash}p{\dimexpr0.68\linewidth-2\tabcolsep\relax}>{\raggedright\arraybackslash}p{\dimexpr0.23\linewidth-2\tabcolsep\relax}@{}}
\caption{Proposed mission requirements for the EPS engineering enterprise.}\label{tab:framework-mission-requirements}\\
\toprule
\textbf{ID} & \textbf{Enterprise mission requirement} & \textbf{Motivating sources/phases} \\
\midrule
\endfirsthead
\toprule
\textbf{ID} & \textbf{Enterprise mission requirement} & \textbf{Motivating sources/phases} \\
\midrule
\endhead
\bottomrule
\endfoot
MR-01 & The EPS engineering enterprise shall obtain MCR acceptance at the end of Pre-Phase A for a mission concept whose EPS feasibility is justified against agreed stakeholder expectations and mission constraints, supported by initial ConOps, budgets, assumptions, risks and a preliminary verification and validation approach. & Pre-Phase A, concept studies \citep[Sec.~3.3]{nasa2016systems}. \\
MR-02 & The enterprise shall obtain SRR acceptance during Phase A for validated EPS requirements and interfaces traced to the mission concept and allocated spacecraft requirements, with measurable criteria consistent with stakeholder expectations. It shall derive, refine and validate EPS and component requirements and interfaces as the design matures through Phase B and adapted C. & A, refinement through B and adapted C \citep[Secs.~3.4--3.6 and~4.1--4.2]{nasa2016systems}. \\
MR-03 & The enterprise shall obtain MDR/SDR acceptance at the end of Phase A for a credible EPS architecture with requirement allocations and justified alternatives, and PDR acceptance at the end of Phase B for a sized, compatible preliminary design supported by sizing and trades, margins, risk dispositions, controlled interfaces and verification and validation plans against applicable requirements and agreed resource and risk limits. & A/B, architecture and preliminary design \citep[Secs.~3.4--3.5]{nasa2016systems}. \\
MR-04 & The enterprise shall obtain CDR acceptance during adapted Phase C, before deliverable coding, for complete detailed EPS model/software design, interface specifications and implementation, integration and verification and validation plans, followed by SIR acceptance at the end of C for released models and declared host-executed control/interface software, applicable MIL/SIL evidence, compatible interfaces and ready integration configurations and procedures. & Adapted C \citep[Sec.~3.6]{nasa2016systems}. Scoped model/software assessment \citep[Sec.~2]{review-0573} \citep[pp.~7--9]{handpicked-0001}. \\
MR-05 & The enterprise shall establish and maintain coherent phase baselines through authorized changes, repeated affected checks and applicable verification and validation evidence, retaining configuration-linked requirements, expectations, methods, results, limitations and finding dispositions. It shall obtain and retain acceptance decisions from the designated review authority against agreed success criteria, prevent progression with blocking findings, and record conditional limits, closure actions, owners and dates. & All scoped phases, technical management \citep[Secs.~6.2--6.7]{nasa2016systems}. \\
\end{longtable}
}

{\small
\setlength{\tabcolsep}{4pt}
\renewcommand{\arraystretch}{1.08}
\begin{longtable}{@{}>{\raggedright\arraybackslash}p{0.12\linewidth}>{\raggedright\arraybackslash}p{\dimexpr0.88\linewidth-2\tabcolsep\relax}@{}}
\caption{Crosscutting enterprise measures of effectiveness.}\label{tab:framework-lifecycle-moes}\\
\toprule
\textbf{ID} & \textbf{Measure and observation basis} \\
\midrule
\endfirsthead
\toprule
\textbf{ID} & \textbf{Measure and observation basis} \\
\midrule
\endhead
\bottomrule
\endfoot
MOE-01 & Engineering burden: active person-hours and elapsed hours for correct requirements, design, change and review tasks, with waiting and unattended execution reported separately. \\
MOE-02 & Early issue detection: expert-confirmed issues found before scoped final design, critical misses and unnecessary investigations against a declared reference issue set. \\
MOE-03 & Evidence and review usefulness: applicable requirement and expectation coverage, critical review defects, reviewer interpretation accuracy and successful independent reconstruction. \\
MOE-04 & Adoption and affordability: person-hours to learn and complete correct tasks and adapt to a mission, and direct framework usage cost in CAD over a declared period and price date. \\
\end{longtable}
}

{\small
\setlength{\tabcolsep}{4pt}
\renewcommand{\arraystretch}{1.08}
\begin{longtable}{@{}>{\raggedright\arraybackslash}p{0.12\linewidth}>{\raggedright\arraybackslash}p{\dimexpr0.33\linewidth-2\tabcolsep\relax}>{\raggedright\arraybackslash}p{\dimexpr0.55\linewidth-2\tabcolsep\relax}@{}}
\caption{Primary root concern contributions motivating the MOEs.}\label{tab:framework-moe-concerns}\\
\toprule
\textbf{MOE} & \textbf{Primary root concerns} & \textbf{Motivation} \\
\midrule
\endfirsthead
\toprule
\textbf{MOE} & \textbf{Primary root concerns} & \textbf{Motivation} \\
\midrule
\endhead
\bottomrule
\endfoot
MOE-01 & \ConcernRef{C-PM-01}, \ConcernRef{C-PM-02}, \ConcernRef{C-SE-03} & Requirements effort, budgets and subsystem changes. \\
MOE-02 & \ConcernRef{C-PM-03}, \ConcernRef{C-EPS-01} & Late redesign and repeated verification burden. \\
MOE-03 & \ConcernRef{C-SE-02}, \ConcernRef{C-EPS-11} & Communicating evidence and designs during reviews. \\
MOE-04 & \ConcernRef{C-PM-01}, \ConcernRef{C-PM-05} & Learning and tool costs under C-PM-01, and mission adaptation. \\
\end{longtable}
}
